% Drawing 138/01-1 -- folder 138, batch 01, archive page 18 (PDF page 19).
% Transcription: une ovale sur laquelle sont marqués les points $0$, $1$, $\infty$
% Kind: oval with three points. Method 2 (schematic redraw from the facsimile). Status: ok.
% Notes: An oval (the real projective line) with the points infinity, 1, 0. A small thickening of the stroke on the upper left is not a marked point (no label). The text below belongs to the transcription.
% Coordinates: 1 unit = 1 mm = 1/100 of the crop width (bbox in page fractions: [0.18, 0.58, 0.42, 0.68]).
% Redrawn by Claude (claude-opus-5-5), issue #73, 2026-11-07; not checked by a human; not compiled (no TeX in the session).
% Rights: reproduces part of the work -- behind RENDER_TRANSCRIPTS, never published by itself.
\documentclass[tikz,border=3pt]{standalone}
\usepackage{amsmath,amssymb}
\usetikzlibrary{patterns,arrows.meta}
\tikzset{pencil/.style={black!45}, struck/.style={black!35}, mark/.style={circle,draw,inner sep=1.3pt}, pt/.style={circle,fill,inner sep=0.9pt}, opt/.style={circle,draw,fill=white,inner sep=1.1pt}}
\begin{document}
\begin{tikzpicture}[x=1mm,y=1mm,line cap=round,line join=round,font=\scriptsize]
\draw plot[smooth cycle, tension=0.55] coordinates {(23,-30) (28,-24) (43,-20) (60,-19) (74,-24) (78,-31) (70,-41) (50,-43) (37,-42) (27,-37)};
\node[pt] at (43,-20.5) {};
\node[pt] at (78,-31) {};
\node[pt] at (37,-42) {};
\node[inner sep=1pt] at (39,-16.5) {\scriptsize $\infty$};
\node[inner sep=1pt] at (82.5,-33) {\scriptsize $1$};
\node[inner sep=1pt] at (37,-47) {\scriptsize $0$};
\end{tikzpicture}
\end{document}
